\section{本章小结}
\label{sec:info-summary}

章首问了三件事：一个随机结果平均含有多少不确定性，观测另一个变量能减少多少，以及信息不够时辨别会受到怎样的限制。本章的回答都建立在自信息$-\ln p$上。它由“罕见更意外、独立信息相加”两条要求在相差一个单位常数的意义下唯一确定，以nat或bit计只是刻度不同。

熵是自信息的平均，等于按最佳方式平均要问的是非题数（以bit计，最佳方案与熵相差不到一题）。有限字母表上熵在$0$与$\ln m$之间，两端分别对应确定结果和均匀分布，并且是概率表的凹函数；可数字母表上熵可以为无穷。对密度计算的微分熵可以为负，随坐标伸缩改变$\ln\lvert a\rvert$，它是量化熵减去精度代价$-\ln\Delta$后剩下的部分，二者的极限关系另需正则性与尾部条件。

两个变量时，联合熵不由边缘决定，条件熵是逐条件值的剩余不确定性的平均，链式法则把整体不确定性按揭示顺序分成几段。互信息定义为联合分布相对边缘乘积的KL散度，非负、对称，为零当且仅当独立，可以为无穷，也不要求密度；熵有限时它等于熵之差，由此得到条件化平均不增加熵与联合熵的次可加性。条件互信息衡量已知第三个变量后的共享信息，它与互信息之间没有固定的大小关系。

表~\ref{tab:info-summary}把本章各个量与它们回答的问题、成立条件对照列出。

\begin{BookInlineTable}
  \centering
  \caption{本章的信息量与主要条件}
  \label{tab:info-summary}
  \begin{tabular}{@{}lll@{}}
    \toprule
    量或结论 & 回答的问题 & 需要核对的条件\\
    \midrule
    $H(X)$ & 平均不确定性 & 可数时可能为无穷\\
    $h(X)$ & 密度的平均对数代价 & 可积性；依赖坐标\\
    $H(X\mid Y)$ & 观测后的剩余不确定性 & 只对平均成立\\
    $I(X;Y)$ & 共享信息 & 熵差需各项有限\\
    $H(P,Q)$ & 按参考分布付出的代价 & 零参考概率；$H(P)<\infty$\\
    数据处理 & 后处理不增加信息 & $X\perp\!\!\!\perp Z\mid Y$\\
    Fano & 辨别错误的下界 & 有限标签与猜错标准\\
    InfoNCE & 互信息的可计算下界 & 不超过$\ln N$；总体值\\
    连续性 & 分布接近时信息量接近 & 有限字母表\\
    \bottomrule
  \end{tabular}
\end{BookInlineTable}

交叉熵把平均代价分成固有部分$H(P)$与失配部分$D_{\mathrm{KL}}(P\Vert Q)$；连续时有确定意义的是后者，熵为无穷时也不能用两者之差定义KL。数据处理不等式由互信息链式法则的两种展开得到，前提是处理中的额外随机性在给定观测时与原对象独立；损失恰是条件互信息，可逆变换不损失信息，但不可逆的处理也可能不损失，判据是反向的条件独立。

信息的上界给出辨别的下界：Fano不等式用剩余不确定性限制有限标签的猜错概率，标签均匀时改写为互信息形式，成为通信中超容量不可靠、统计中候选太密无法区分的共同工具。反方向的InfoNCE用辨认真配对的总体损失给出互信息的下界，但下界受候选数$\ln N$与评分函数类的限制，下界小不能推出独立，单批样本值也不自动是下界。最后，有限字母表上熵与互信息对总变差连续，界中含有$\theta\ln(1/\theta)$与字母表大小；无限字母表或连续空间中，接近的分布可以有相差很远的信息量，弱收敛本身不保证熵收敛。

后续各章分别使用这些量。第~\ref{chap:lossless-coding}章到第~\ref{chap:rate-distortion}章研究用熵、互信息衡量的最短描述与有失真描述；第~\ref{chap:discrete-channel-capacity}章到第~\ref{chap:channel-codes}章把容量定义为互信息的最大值，并用Fano不等式证明逆定理；第~\ref{chap:dependence-measures}章把互信息与基于距离、核和离散化的依赖度量放在一起比较。
